\section{Introduction}

The documents in DocVQA-2026 are long, multi-page, high-resolution, and largely visual. A single question can range over dozens of pages of figures, plots, schematics, maps, and infographics, and the answer-bearing content is frequently non-textual --- a value in a chart cell, a label on an engineering drawing, a region of a poster. Because much of the signal never resolves to text, OCR-and-text pipelines cannot reach it. Answering is, first, a problem of \emph{finding}: locating the right page and the right region before any reading can begin, then often combining or computing over what is found.

Scaling the model does not by itself crack this. On the DocVQA-2026 validation split, bare setups score low --- a raw multi-image VLM call, a multimodal agent reading pages in its own context, and a competition-kit prompt with no scaffold all land near 17--22 ANLS. Frontier-scale frozen models do better but only plateau: the official ICDAR-2026 baselines place Gemini~3~Pro at 37.5 and Gemini~3~Flash at 33.75 ANLS on the same validation split, matching or modestly exceeding the bare agents but low for their size on documents this large. Both ends of the scale axis fall short of what the task affords, which points away from model capacity as the binding constraint.

What moves accuracy is \emph{active perception}. Give a code-capable reasoner a persistent Python REPL and a single visual primitive --- an on-demand call to a frozen VLM perception tool against an arbitrary image region --- and the reasoner can direct perception rather than consume it whole: crop to the evidence, zoom for acuity, and do in code the coordinate and numeric arithmetic the VLM cannot. Holding both the reasoner and the frozen VLM fixed and varying only this scaffold spans roughly 25 ANLS points. The REPL active-perception scaffolds reach $\sim$39 ANLS --- far above tool-based ReAct, which has the same VLM but no REPL ($\sim$25), and above no-REPL VLM agents ($\sim$20) --- and match or exceed the much larger frontier frozen models above. The result is demonstrated chiefly at a Qwen3.5-27B backbone, but it is not specific to one model: it holds across model families and is gated by base capability rather than size.

A second thread follows from the structure of the winning scaffold. Active perception is a \emph{family} of harnesses, and its append-only member preserves a growing-prefix trajectory --- the structure that weight-level training assumes --- at no measured accuracy cost, making it the principled target for fine-tuning. In a preliminary study, we ask whether training a small reasoner (Qwen3.5-4B) inside that member can lift a $\leq$8B agent further toward the gain.\footnote{The original proposal envisioned a fuller progression --- supervised fine-tuning, then reinforcement learning, then on-policy distillation. We report the supervised stage and leave the on-policy methods to future work. Agentic reinforcement learning in this setting is bottlenecked by on-policy rollout generation --- each update needs many long, multi-turn trajectories, and every perception step queries a slow frozen vision--language model --- and the hybrid-attention backbone added integration overhead newer to the open training stack; within the project's time budget, the controlled harness evaluation (\S\ref{sec:harness}) was the more valuable investment.}

We make the following contributions.\footnote{Code is available at \url{https://github.com/bdsaglam/docvqa} (evaluation) and \url{https://github.com/bdsaglam/docvqa-verl} (fine-tuning).}

\begin{enumerate}
\item \textbf{Active perception is the dominant lever on document-agent accuracy.} A controlled comparison on DocVQA-2026 val (binary ANLS at threshold 0.9, $n{=}8$, frozen Qwen3.5-27B VLM) shows REPL active-perception scaffolds ($\sim$39 ANLS) far surpassing tool-based ReAct ($\sim$25) and no-REPL VLM agents ($\sim$20), and matching or exceeding the official ICDAR-2026 frontier baselines on the same validation split (Gemini~3~Pro 37.5, Gemini~3~Flash 33.75 --- low for their size on documents this large). The ordering is robust across two model families (Qwen3.5, Gemma-4) and four reasoner sizes (\S\ref{sec:harness}): the lever is the scaffold, not the model.

\item \textbf{Why and when active perception works.} The code REPL and visual perception are each essential --- dropping either collapses the score --- and the advantage concentrates on \emph{visually dense} pages, where cropping recovers detail a whole-page read misses (the lead over ReAct is largest on the detail-dense categories: engineering drawings, business reports, maps; \S\ref{sec:mechanism}); OCR-free visual perception is decisive (swapping vision for OCR text costs $-25.5$ points, with engineering drawings and maps falling to zero). The mechanism is that the REPL converts reasoning capability into targeted perception --- which is also why the advantage is \emph{capability-gated}, not scale-gated: a modern Qwen3.5-4B clears the bar, while a same-size Gemma-4-E4B (no lift) and an older, larger Qwen3-8B (where the REPL actively hurts a weak coder) are clean negative controls.

\item \textbf{A preliminary $\leq$8B training study.} We build a rejection-sampling fine-tuning pipeline and a training-data pool for the trainable scaffold and run a preliminary LoRA study on Qwen3.5-4B; reinforcement learning and on-policy distillation are motivated but left to future work, and the quantitative verdict awaits a corrected evaluation.
\end{enumerate}

The remainder proceeds from background and task (\S\ref{sec:background}) to the active-perception evaluation and its mechanism (\S\ref{sec:harness}), the argument for the append-only scaffold as a trainable target (\S\ref{sec:scaffold}), the preliminary small-agent training study (\S\ref{sec:training}), related work (\S\ref{sec:related}), and conclusions (\S\ref{sec:conclusion}).
